% Generated draft: requires mathematical and source review.
\documentclass[11pt]{article}
\usepackage[margin=25mm]{geometry}
\usepackage{fontspec}
\setmainfont{Latin Modern Roman}
\usepackage{amsmath,amssymb,mathtools,unicode-math}
\setmathfont{Latin Modern Math}
\usepackage{enumitem,needspace,xurl,hyperref}
\setlength{\parindent}{0pt}
\setlength{\parskip}{5pt}
\setlength{\emergencystretch}{3em}
\newcommand{\N}{\mathbb{N}}
\newcommand{\Z}{\mathbb{Z}}
\newcommand{\Q}{\mathbb{Q}}
\newcommand{\R}{\mathbb{R}}
\newcommand{\C}{\mathbb{C}}
\newcommand{\F}{\mathbb{F}}
\newcommand{\A}{\mathbb{A}}
\newcommand{\PP}{\mathbb{P}}
\newcommand{\E}{\mathbb{E}}
\newcommand{\eps}{\varepsilon}
\newcommand{\dd}{\,\mathrm{d}}
\begin{document}
\section*{UnsolvedMath ID: 30001874}
\textbf{Source problem code:} catalogue entry
\subsection{Definitions and mathematical statement}
The supplied record asks for a characterization of nonnegative functions \(F\) for which a trading strategy can guarantee
\[
X_t^*=\max_{s\le t}X_s,
\qquad
\mathrm{IK}_t\ge F(X_t^*,X_t)
\quad\text{for all }t,
\]
while making \(F(X_t^*,X_t)\) grow ``as nearly as possible like'' \(X_t^*\). Here \(X_t^*\) denotes the running maximum appearing in the source, whenever that maximum exists. The expression \(\mathrm{IK}_t\) preserves the supplied token \texttt{IK\_t}; its meaning and intended typography are not established by the record.

This is an incomplete characterization request, not a fully specified mathematical proposition. The record does not specify the time index or its initial boundary, the domain and admissible trajectories of \(X\), hypotheses ensuring that the maximum is attained, the meaning of \(\mathrm{IK}_t\), initial capital, trading rules, or admissible strategies. It also does not specify the domain or regularity of \(F\), beyond nonnegativity, or whether the guarantee is pathwise, almost sure, or interpreted otherwise. The phrase describing growth supplies no asymptotic regime, comparison criterion, or optimization objective. Accordingly, these missing conditions are left unresolved; the displayed expression is only the supplied target schema.

The registry labels the entry open, but its supplied statement-verification note says that the source immediately answers the characterization request in Theorems 4 and 5 and that this is an answered setup rather than an unresolved problem. Those theorem statements are not supplied. Thus this draft does not assert an open conjecture or reproduce a characterization theorem.
\subsection{Short English statement}
The source asks which nonnegative functions of a process and its running maximum can be guaranteed by trading, with growth close to that maximum. Essential market assumptions and the growth criterion are missing. Supplied verification describes the question as already answered in the source, despite the registry's open label; human review is required.
\subsection{Sources}
\begin{itemize}
\item Supplied local UnsolvedMath registry record, permanent ID \(30001874\), exact source problem code \texttt{catalogue entry}, titled ``Answered Characterization of Lookback Adjusters,'' in the Oberwolfach Reports Open Problems collection. No proposer is identified; the supplied proposed year is 2011.
\item Supplied report citation: \emph{Mini-Workshop: Mathematics of Machine Learning (2012)}, Oberwolfach Reports, DOI \texttt{10.4171/owr/2011/42}, pp.\ 2365--2396. Supplied source URL: \url{https://doi.org/10.4171/owr/2011/42}. The background attributes the material to Oberwolfach workshop reports and their open-problem discussions. The report contents were not independently consulted for this draft.
\item Supplied literature source, labeled ``original OWR report'': \url{https://doi.org/10.4171/owr/2011/42}.
\item Literature assessment dated exactly \texttt{2026-08-21}: supplied dated curation, not independently verified. Its status label is open, and its assessment describes the extracted text as a qualitative discussion of trading-risk tradeoffs rather than a definite theorem or claim with a verifiable resolution. The background calls this dated literature triage and requests independent verification of time-sensitive status and cited claims.
\item Supplied statement status: \texttt{unrecoverable}. The statement-verification note says that necessary-and-sufficient characterizations appear immediately in Theorems 4 and 5, calls the extract an answered setup, and states that no separate clean statement is promoted; the original remains publicly visible for provenance. These are supplied source-verification claims, not independently verified mathematical conclusions.
\end{itemize}
\end{document}
